\section{Why Formal Verification}


Software testing can show the presence of bugs, never their absence. For blockchain systems
where bugs cause irreversible fund loss, this is insufficient. Formal verification uses
mathematical proof to establish that a program or protocol satisfies its specification for
\emph{all} inputs, not just the ones a test suite happens to exercise.

\subsection{The Cost of Bugs in Blockchain}

The DAO hack (2016, \$60M), Wormhole bridge exploit (2022, \$320M), and Ronin bridge
compromise (2022, \$625M) share a common pattern: the vulnerable code was tested, audited,
and deployed by competent teams. Testing and auditing are necessary but not sufficient.
Formal verification closes the gap by providing machine-checked guarantees.

\subsection{Why Not Just Test More?}

Consider a BFT consensus protocol with $n$ validators. Testing safety requires exploring
an exponential state space: $2^n$ possible failure combinations, arbitrary message
reorderings, unbounded network delays. Even property-based testing with millions of random
inputs covers a negligible fraction of this space. Formal proof covers it all by reasoning
about the structure of the state space rather than sampling points within it.

\subsection{Proof Assistants vs.\ Pen-and-Paper Proofs}

Peer-reviewed papers in distributed systems regularly contain subtle errors in their proofs.
Hawblitzel et al.\ found errors in published proofs of Paxos variants. A proof assistant
like Lean~4 mechanically checks every step: if the proof compiles, the theorem holds
(modulo the axioms of the underlying logic, which for Lean~4 is the Calculus of Inductive
Constructions plus propositional extensionality and quotient types).

%% ============================================================
